\SourceLeaf{65}
\begin{center}
\textit{SIXTH PROOF OF THE FUNDAMENTAL THEOREM}\par
\textit{IN THE THEORY OF QUADRATIC RESIDUES.}
\end{center}

\Article{1}

\textsc{Theorem.}\quad \textit{Let $p$ denote a prime number (positive and odd), $n$ a positive integer not divisible by $p$, and $x$ an indeterminate quantity. Then the function}

\[
\begin{aligned}
1+x^n+x^{2n}+x^{3n}+\text{etc.}+x^{np-n}
\end{aligned}
\]

\textit{is divisible by}

\[
\begin{aligned}
1+x+xx+x^3+\text{etc.}+x^{p-1}
\end{aligned}
\]

\textsc{Proof.}\quad Take a positive integer $g$ such that $gn\equiv1\quad(\text{mod. }p)$, and put $gn=1+hp$.\quad Then we have

{\small
\[
\begin{aligned}
\frac{1+x^n+x^{2n}+x^{3n}+\text{etc.}\,x^{np-n}}{1+x+xx+x^3+\text{etc.}\,x^{p-1}} &=\frac{(1-x^{np})(1-x)}{(1-x^n)(1-x^p)}\\&=\frac{(1-x^{np})(1-x^{gn}-x+x^{hp+1})}{(1-x^n)(1-x^p)}\\
&=\frac{1-x^{np}}{1-x^p}\cdot\frac{1-x^{gn}}{1-x^n}-\frac{x(1-x^{np})}{1-x^n}\cdot\frac{1-x^{hp}}{1-x^p}
\end{aligned}
\]
}

which is evidently a polynomial.\quad Q. E. D.

Thus every polynomial in $x$ divisible by $\frac{1-x^{np}}{1-x^n}$ is also divisible by $\frac{1-x^p}{1-x}$.

\Article{2}

Let $\alpha$ denote a positive primitive root modulo $p$; that is, let $\alpha$ be a positive integer such that the least positive residues of the powers $1,\alpha,\alpha\alpha,\alpha^3\dotsfive\alpha^{p-2}$ modulo $p$ are precisely the numbers $1,2,3,4\dotsfive p-1$, without regard to order.\quad Further, denoting by $fx$ the function

\[
\begin{aligned}
x+x^\alpha+x^{\alpha\alpha}+x^{\alpha^3}+\text{etc.}+x^{\alpha^{p-2}}+1
\end{aligned}
\]

it is clear that $fx-1- x-xx-x^3-\text{etc.}-x^{p-1}$ is divisible by $1-x^p$, and therefore all the more by $\frac{1-x^p}{1-x}=1+x+xx+x^3+\text{etc.}+x^{p-1}$. This last function therefore also divides $fx$ itself.\quad Since $x$ represents an indeterminate quantity, it follows that $f(x^n)$ is also divisible by $\frac{1-x^{np}}{1-x^n}$, and hence, by the preceding article, by $\frac{1-x^p}{1-x}$, whenever $n$ is an integer not divisible by $p$. If, on the other hand, $n$ is an integer divisible by $p$, the individual terms of the function $f(x^n)$, each decreased by\SourceLeaf{66} one, are divisible by $1-x^p$. In this case, therefore, $f(x^n)-p$ is also divisible by $1-x^p$, and hence by $\frac{1-x^p}{1-x}$.

\Article{3}

\textsc{Theorem.}\quad \textit{Putting}

\[
\begin{aligned}
x-x^\alpha+x^{\alpha\alpha}-x^{\alpha^3}+x^{\alpha^4}-\text{etc.}-x^{\alpha^{p-2}}=\xi
\end{aligned}
\]

\textit{then $\xi\xi\mp p$ is divisible by $\frac{1-x^p}{1-x}$, taking the upper sign whenever $p$ is of the form $4k+1$, and the lower whenever $p$ is of the form $4k+3$.}

\textsc{Proof.}\quad It is easy to see that, of the following $p-1$ functions

\[
\begin{aligned}
&+x\xi-xx+x^{\alpha+1}-x^{\alpha\alpha+1}\\&\qquad+\text{etc.}+x^{\alpha^{p-2}+1}\\
&-x^\alpha\xi-x^{2\alpha}+x^{\alpha\alpha+\alpha}-x^{\alpha^3+\alpha}\\&\qquad+\text{etc.}+x^{\alpha^{p-1}+\alpha}\\
&+x^{\alpha\alpha}\xi-x^{2\alpha\alpha}+x^{\alpha^3+\alpha\alpha}-x^{\alpha^4+\alpha\alpha}\\&\qquad+\text{etc.}+x^{\alpha^p+\alpha\alpha}\\
&-x^{\alpha^3}\xi-x^{2\alpha^3}+x^{\alpha^4+\alpha^3}-x^{\alpha^5+\alpha^3}\\&\qquad+\text{etc.}+x^{\alpha^{p+1}+\alpha^3}
\end{aligned}
\]

etc., as far as

\[
\begin{aligned}
-x^{\alpha^{p-2}}\xi-x^{2\alpha^{p-2}}+x^{\alpha^{p-1}+\alpha^{p-2}}-x^{\alpha^p+\alpha^{p-2}}\\&\qquad+\text{etc.}+x^{\alpha^{2p-4}+\alpha^{p-2}}
\end{aligned}
\]

the first is $=0$, while each of the others is divisible by $1-x^p$. Their sum is therefore also divisible by $1-x^p$, and is found to be

\[
\begin{aligned}
&=\xi\xi-(f(xx)-1)+(f(x^{\alpha+1})-1)-(f(x^{\alpha\alpha+1})-1)\\&\qquad+(f(x^{\alpha^3+1})-1)-\text{etc.}\\
&\qquad+(f(x^{\alpha^{p-2}+1})-1)\\
&=\xi\xi-f(xx)+f(x^{\alpha+1})-f(x^{\alpha\alpha+1})\\&\qquad+f(x^{\alpha^3+1})-\text{etc.}+f(x^{\alpha^{p-2}+1})=\Omega
\end{aligned}
\]

\Needspace{13\baselineskip}
Thus this expression $\Omega$ is also divisible by $\frac{1-x^p}{1-x}$. Now, among the exponents $2,\ \alpha+1,\ \alpha\alpha+1,\ \alpha^3+1\dotsfive\alpha^{p-2}+1$, only one is divisible by $p$, namely $\alpha^{\frac12(p-1)}+1$. Hence, by the preceding article, the following individual parts of the expression $\Omega$

\[
\begin{aligned}
f(xx),\quad f(x^{\alpha+1}),\quad f(x^{\alpha\alpha+1}),\quad (fx^{\alpha^3+1})\quad\text{etc.}
\end{aligned}
\]

with the sole exception of the term $f(x^{\alpha^{\frac12(p-1)}+1})$, are divisible by $\frac{1-x^p}{1-x}$. Those parts may therefore be removed, leaving the following function still divisible by $\frac{1-x^p}{1-x}$:

\[
\begin{aligned}
\xi\xi\mp f(x^{\alpha^{\frac12(p-1)}+1})
\end{aligned}
\]

\SourceLeaf{67} Here the upper or lower sign applies according as $p$ is of the form $4k+1$ or of the form $4k+3$. Since, moreover, $f(x^{\alpha^{\frac12(p-1)}+1})-p$ is divisible by $\frac{1-x^p}{1-x}$, $\xi\xi\mp p$ is also divisible by $\frac{1-x^p}{1-x}$. Q. E. D.

To prevent any ambiguity from the double sign, we shall use $\varepsilon$ to denote the number $+1$ or $-1$, according as $p$ is of the form $4k+1$ or $4k+3$. Thus $\frac{(1-x)(\xi\xi-\varepsilon p)}{1-x^p}$ is a polynomial in $x$, which we shall denote by $Z$.

\Article{4}

Let $q$ be a positive odd number, so that $\frac12(q-1)$ is an integer. Then $(\xi\xi)^{\frac12(q-1)}-(\varepsilon p)^{\frac12(q-1)}$ is divisible by $\xi\xi-\varepsilon p$, and hence also by $\frac{1-x^p}{1-x}$. Let us put $\varepsilon^{\frac12(q-1)}=\delta$, and

\[
\begin{aligned}
\xi^{q-1}-\delta p^{\frac12(q-1)}=\frac{1-x^p}{1-x}\cdot Y
\end{aligned}
\]

Then $Y$ is a polynomial in $x$, and $\delta=\mathord{.}+1$ whenever one of the numbers $p,q$, or indeed both, is of the form $4k+1$. On the other hand, $\delta=-1$ whenever both $p,q$ are of the form $4k+3$.

\newpage
\Article{5}

Let us now suppose that $q$ too is a prime number, different from $p$. It is clear from the theorem proved in Article 51 of the \textit{Disquisitiones Arithmeticae} that

\[
\begin{aligned}
\mathord{.}\quad\xi^q-(x^q-x^{q\alpha}+x^{q\alpha\alpha}-x^{q\alpha^3}+\text{etc.}-x^{q\alpha^{p-2}})
\end{aligned}
\]

is divisible by $q$, or is of the form $qX$, where $X$ is a polynomial in $x$ whose numerical coefficients are also integers (as is also to be understood of the other polynomials $Z,Y,W$ occurring here). For the modulus $p$ and primitive root $\alpha$, let us denote the index of the number $q$ by $\mu$; that is, let $q\equiv\alpha^\mu\ (\mathrm{mod.}\ p)$. The numbers $q,\ q\alpha,\ q\alpha\alpha,\ q\alpha^3\dotsfive q\alpha^{p-2}$ are therefore congruent modulo $p$, respectively, to the numbers $\alpha^\mu,\allowbreak\ \alpha^{\mu+1},\allowbreak\ \alpha^{\mu+2}\dotsfour\alpha^{p-2},\allowbreak\ 1,\allowbreak\ \alpha,\allowbreak\ \alpha\alpha\dotsfour\alpha^{\mu-1}$, so that

\[
\begin{aligned}
&x^q-x^{\alpha^\mu}\\
&x^{q\alpha}-x^{\alpha^{1+1}}\\
&x^{q\alpha\alpha}-x^{\alpha^{\mu+2}}\\
&x^{q\alpha^3}-x^{\alpha^{\mu+3}}\\
&\hspace{9mm}\begin{gathered}.\\.\\.\\.\end{gathered}
\end{aligned}
\]

\SourceLeaf{68}
\[
\begin{aligned}
&x^{q\alpha^{p-\mu-2}}-x^{\alpha^{p-2}}\\
&x^{q\alpha^{p-\mu-1}}-x\\
&x^{q\alpha^{p-\mu}}-x^\alpha\\
&x^{q\alpha^{p-\mu+1}}-x^{\alpha\alpha}\\
&\hspace{9mm}\begin{gathered}.\\.\\.\\.\end{gathered}\\
&x^{q\alpha^{p-2}}-x^{\alpha^{\mu-1}}
\end{aligned}
\]

\Needspace{14\baselineskip}
are divisible by $1-x^p$. Taking these quantities alternately with positive and negative signs and adding them shows that $1-x^p$ divides the function

\[
\begin{aligned}
x^q-x^{q\alpha}+x^{q\alpha\alpha}-x^{q\alpha^3}+\text{etc.}-x^{q\alpha^{p-2}}\mp\xi
\end{aligned}
\]

with the upper or lower sign applying according as $\mu$ is even or odd; that is, according as $q$ is a quadratic residue or a nonresidue of $p$. We shall therefore put

\[
\begin{aligned}
x^q-x^{q\alpha}+x^{q\alpha\alpha}- x^{q\alpha^3}+\text{etc.}-x^{q\alpha^{p-2}}-\gamma\xi\\&=(1-x^p)W
\end{aligned}
\]

taking $\gamma=+1$ or $\gamma=-1$, according as $q$ is a quadratic residue or a nonresidue of $p$. It is clear that $W$ is a polynomial.

\Needspace{13\baselineskip}
\Article{6}

With these preparations, we combine the preceding equations to obtain

\[
\begin{aligned}
q\xi X=\varepsilon p(\delta p^{\frac12(q-1)}-\gamma)\\&\qquad+\frac{1-x^p}{1-x}\cdot(Z(\delta p^{\frac12(q-1)}-\gamma)\\&\qquad+Y\xi\xi-W\xi(1-x))
\end{aligned}
\]

Let us suppose that dividing the function $\xi X$ by

\[
\begin{aligned}
x^{p-1}+x^{p-2}+x^{p-3}+\text{etc.}+x+1
\end{aligned}
\]

yields the quotient $U$ with remainder $T$, or that we have

\[
\begin{aligned}
\xi X=\frac{1-x^p}{1-x}\cdot U+T
\end{aligned}
\]

where $U,T$ are polynomials whose numerical coefficients are also integers, and $T$ is certainly of lower degree than the divisor. We therefore have

\[
\begin{aligned}
qT-\varepsilon p(\delta p^{\frac12(q-1)}-\gamma)\\&=\frac{1-x^p}{1-x}\cdot(Z(\delta p^{\frac12(q-1)}-\gamma)\\&\qquad+Y\xi\xi-W\xi(1-x)-qU)
\end{aligned}
\]

This equation evidently cannot hold unless the left side and the right side each vanish separately. Thus $\varepsilon p(\delta p^{\frac12(q-1)}-\gamma)$ is divisible by $q$\SourceLeaf{69}, as is $\delta p^{\frac12(q-1)}-\gamma$; and therefore, since $\delta\delta=1$, the number $p^{\frac12(q-1)}-\gamma\delta$ is also divisible by $q$.

\Needspace{8\baselineskip}
Now, if $\beta$ denotes unity taken positively or negatively, according as $p$ is a quadratic residue or a quadratic nonresidue of the number $q$, then $p^{\frac12(q-1)}-\beta$ is divisible by $q$, and hence so is $\beta-\gamma\delta$. This cannot happen unless $\beta=\gamma\delta$. The fundamental theorem follows immediately. Namely,

I. Whenever the numbers $p,q$ both have the form $4k+1$, or just one of them does, so that $\delta=+1$, we have $\beta=\gamma$. Consequently either $q$ is a quadratic residue of $p$ and at the same time $p$ a quadratic residue of $q$; or $q$ is a nonresidue of $p$ and at the same time $p$ a nonresidue of $q$.

II. Whenever both $p,q$ are of the form $4k+3$, so that $\delta=-1$, we have $\beta=-\gamma$. Consequently either $q$ is a quadratic residue of $p$ and at the same time $p$ a nonresidue of $q$; or $q$ is a nonresidue of $p$ and at the same time $p$ a residue of $q$.

Q. E. D.
